\documentclass[10pt,a4paper]{amsart}
\usepackage[margin=19mm]{geometry}
\usepackage{amsmath,amssymb,mathtools}
\usepackage[scr=rsfs,cal=euler]{mathalfa}
\usepackage{xfrac}
\usepackage[unicode,hidelinks,bookmarks=false]{hyperref}
\newcommand{\ie}{i.e.\ }
\newcommand{\cf}{cf.\ }
\newcommand{\resp}{respectively\ }
\newcommand{\Subsection}{\subsection}
\providecommand{\nobreakdash}{\nobreak\hbox{-}}
\setlength{\emergencystretch}{2em}
\multlinegap0pt
\usepackage{amssymb}
\newtheorem{example}{Example}
\newtheorem{lemm}{Lemma}
\newtheorem{prop}{Proposition}
\title{Varieties with representable ${\rm CH}_0$-group and a question of Colliot-Th\'{e}l\`{e}ne}
\author{Claire Voisin}
\date{Source: arXiv:2508.02331v2 (CC BY 4.0)}
\begin{document}
\maketitle

\noindent\emph{Source and licence.} Varieties with representable ${\rm CH}_0$-group and a question of Colliot-Th\'{e}l\`{e}ne. Version arXiv:2508.02331v2 (CC BY 4.0), distributed by arXiv under the Creative Commons Attribution 4.0 International licence, http://creativecommons.org/licenses/by/4.0/, as recorded in the arXiv licence metadata for this exact version and shown on the abstract page https://arxiv.org/abs/2508.02331v2. Full licence notice: \texttt{licenses/arxiv-2508.02331-CC-BY-4.0.txt}.\par\smallskip

\noindent\textit{Editorial scope.} Counted unit: the article's prop prokununiv (source0.tex lines 155--163). The article's complete original proof of it is delivered with it (source0.tex lines 168--197, one \texttt{proof} environment of 30 lines, which the article places away from the statement). No mathematics is written, repaired or re-derived: the statement and the proof are the article's own lines, in the article's own notation, and no retained line is substituted or re-flowed.  Retained uncounted as the source-local context the proof needs: lines 63--64; lines 126--152.  Nothing was omitted from any retained span, so the package carries no omission record.  No retained line contains a citation, so the wrapper carries no bibliography and no bibliography entry.  No retained line contains a figure environment, an image-inclusion command or drawing code, so no image is delivered and the package declares no asset dependency.  Editorial formatting only, in addition: an AMS \texttt{amsart} preamble that restates the article's own declarations and macros used by the retained lines, namely the environments \texttt{example}, \texttt{lemm}, \texttt{prop} on separate counters, together with the standard packages those lines need.  The retained environments print in this document as Example 1, Lemm 1, Proposition 1 in order of appearance, whereas the article prints the same items with the numbering of its own full document.  The complete native source of the article as distributed in the arXiv e-print for this exact version is bundled unchanged as \texttt{sources/182660/source0.tex}.
\begin{example}\label{example}{\rm A smooth projective curve admits a universal $0$-cycle given by the Poincar\'{e} divisor on $J(C)\times C $.}
\end{example}

We will use the following construction.
\begin{lemm} \label{lerefsuggestion} Let $X$ be  a smooth projective variety and $C$ a smooth curve with a morphism
$j: C\rightarrow X$ inducing
$$j_*:J(C)\rightarrow {\rm Alb}(X).$$
Assume there exists a morphism of abelian varieties $s:{\rm Alb}(X)\rightarrow J(C)$ such that
\begin{eqnarray} \label{eqjstar1}s\circ j_*= k{\rm Id}_{{\rm Alb}(X)}\end{eqnarray} for some integer $k$. Then there exists a codimension $n$ cycle
$$\Gamma\in{\rm CH}^n({\rm Alb}(X)\times X)$$
such that
$$[\Gamma]_*=k{\rm Id}_{{\rm Alb}(X)}:{\rm Alb}({\rm Alb}(X))={\rm Alb}(X)\rightarrow {\rm Alb}(X).$$
\end{lemm}
\begin{proof}
The curve $C$ admits a universal $0$-cycle $\Gamma_C\in {\rm CH}^1(J(C)\times C)$ (see Example \ref{example}). Then
the codimension $n$-cycle
$$\Gamma_X:=(\overline{s},{\rm Id}_X)^*(({\rm Id}_{J(C)},j)_*(\Gamma_C))\in{\rm CH}^n({\rm Alb}(X)\times X)$$
 has by (\ref{eqjstar1}) the property that
$$ [\Gamma_X]_*=k {\rm Id}_{{\rm Alb}(X)}:{\rm Alb}({\rm Alb}(X))={\rm Alb}(X)\rightarrow {\rm Alb}(X).$$
\end{proof}
As a first application, let us prove that any smooth projective variety $X$ has a universal $0$-cycle with rational coefficients. Indeed, choosing a smooth  curve $C$ which is a complete intersection of ample hypersurfaces in $X$, and denoting $j: C\rightarrow X$ the inclusion map,
$$j_*:H_1(C,\mathbb{Z})\rightarrow H_1(X,\mathbb{Z})_{\rm tf}$$
is a surjective morphism of Hodge structures. By semisimplicity of polarized weight $1$ Hodge structures,
there exists a morphism of Hodge structures
$$ s: H_1({\rm Alb}(X),\mathbb{Z})\cong  H_1(X,\mathbb{Z})_{\rm tf}\rightarrow H_1(C,\mathbb{Z})$$
such that \begin{eqnarray} \label{eqjstar} j_*\circ s=k{\rm Id}_{H_1(X,\mathbb{Z})_{\rm tf}}\end{eqnarray}
for some nonzero integer $k$. The morphism $s$ induces a morphism
$$\overline{s}: {\rm Alb}(X)\rightarrow J(C)$$
of abelian varieties and $j_*,\,\overline{s}$ satisfy the assumptions of Lemma \ref{lerefsuggestion}. Hence Lemma
 \ref{lerefsuggestion} provides a universal $0$-cycle with rational coefficients.

\begin{prop} \label{prokununiv}The following statements are equivalent:

(i)  The K\"{u}nneth projector $\delta_1$ is algebraic on $X\times X$ and is the class of a cycle supported on $X\times C$ for some curve $C\subset X$.

(ii) $X$  has a universal $0$-cycle and the Albanese variety of $X$ is a direct summand in the Jacobian of a curve.

(iii) The motive of $X$ contains a direct summand $M_1$, with the property that $M_1$ is a direct summand in the motive of a curve and ${\rm Alb}(X)={\rm Alb}(M_1)$.

\end{prop}

\begin{proof}[Proof of Proposition \ref{prokununiv}] Assume there exist a  curve $C\subset X$ and a codimension $1$ cycle $Z\subset X\times C$ such that the class of $Z$ in $X$ is $\delta_1$. We can replace by desingularization  $C$ by a smooth curve $\widetilde{C}$ with a morphism $\tilde{j}:\widetilde{C}\rightarrow X$, such that $Z$ lifts to a codimension $1$ cycle $\widetilde{Z}\subset X\times C$. Then, as the class of $Z$ in $X$ is $\delta_{1}$, we have for any $\alpha\in H_1(X,\mathbb{Z})$
\begin{eqnarray}\label{eqpourpreubeproofdirect} \alpha=\delta_{1*}(\alpha)=\tilde{j}_*([\widetilde{Z}]_*(\alpha))\,\,{\rm in}\,\,H_1(X,\mathbb{Z}).
\end{eqnarray}
It follows from (\ref{eqpourpreubeproofdirect}) that, via $\tilde{j}_*:J(\widetilde{C})\rightarrow {\rm Alb}(X)$, ${\rm Alb}(X)$ is a direct summand of $J(\widetilde{C})$, with right  inverse given by $ \sigma:=[\widetilde{Z}]_*:{\rm Alb}(X)\rightarrow J(\widetilde{C})$, which proves the second statement in (ii).
Next, we apply Lemma \ref{lerefsuggestion} to $\tilde{j}$ and $\sigma$ and consider
 the cycle
$$ \Gamma_1:= (\sigma,{\rm Id}_C)^*\circ({\rm Id}_{J(C)},\tilde{j})_*(P_{\widetilde{C}})\in {\rm CH}^n({\rm Alb}(X)\times X ).$$
It has the property that $[\Gamma_{1}]_*:{\rm Alb}({\rm Alb}(X))={\rm Alb}(X)\rightarrow {\rm Alb}(X)$ is the identity, since $\tilde{j}_*\circ \sigma={\rm Id}_{{\rm Alb}(X)}$. Hence $X$ has a universal $0$-cycle. Thus (i) implies (ii).




 In the other direction, let $\Gamma\in {\rm CH}^n( {\rm Alb}(X)\times X)$ be a universal $0$-cycle for $X$. This means that
$$[\Gamma]_*:H_1({\rm Alb}(X),\mathbb{Z})\rightarrow  H_1(X,\mathbb{Z})$$
is the natural isomorphism. Let $D$ be a (nonnecessarily connected) smooth projective  curve such that ${\rm Alb}(X)$   is a direct summand of $J(D)$, and let $$\sigma: {\rm Alb}(X)\rightarrow J(D),\,\pi:J(D)\rightarrow {\rm Alb}(X)$$ be such that $\pi\circ \sigma={\rm Id}_{{\rm Alb}(X)}$. Finally, let $i_D: D\rightarrow J(D)$ be the natural embedding and
consider the cycle
$$\Gamma_1:=    \Gamma\circ \pi\circ i_D\in{\rm CH}^n(D\times X).$$
We have $$[\Gamma_1]_*=[\Gamma]_*\circ \pi_*\circ i_{D*}=\pi: J(D)\rightarrow {\rm Alb}(X).$$
 Next, let $P_D\in {\rm CH}^1( J(D)\times D)$ be a universal $0$-cycle (or Poincar\'{e} divisor) for $D$ and consider next the cycle
$$\Gamma_2=P_D\circ \sigma \circ {\rm alb}_X\in {\rm CH}_n(X\times D).$$
We have
$[\Gamma_2]_*=\sigma:{\rm Alb}(X)\rightarrow J(D)$. It follows that
$\Gamma_X:=\Gamma_1\circ \Gamma_2\in {\rm CH}^n(X\times X)$
has the property that
\begin{eqnarray}\label{eqinversefacteur} [\Gamma_X]_*=[\Gamma_1]_*\circ [\Gamma_2]_*=\pi\circ \sigma={\rm Id}_{\rm Alb}(X).\end{eqnarray}
Finally $\Gamma_X$ is supported on $X\times D'$ where $D'\subset X$ is the curve in $X$ defined as ${\rm pr}_2({\rm Supp}(\Gamma_1)$. Furthermore one easily checks that $[\Gamma_X]^*=0$ on $H^i(X,\mathbb{Z})$ for $i\not =1$. It follows that $[\Gamma_X]=\delta_1$.
This proves that (ii) implies (i).

The equivalence of  (ii) with property (iii) is  now clear. If we have a direct summand $M_1$ as in (iii), then $M_1$ has a universal $0$-cycle, being a direct summand in the motive of a curve $C$, and its Albanese variety is a direct summand of $J(C)$. Hence (ii) holds. Conversely, if (ii) holds, then we construct $\Gamma_1$ and $\Gamma_2$ as above and (\ref{eqinversefacteur}) exactly says that $[\Gamma_2]$ realizes motivically  $H_1(X,\mathbb{Z})$ as a direct summand in $H_1(D,\mathbb{Z})$.
\end{proof}

\end{document}
